% Generated draft: requires mathematical and source review.
\documentclass[11pt]{article}
\usepackage[margin=25mm]{geometry}
\usepackage{fontspec}
\setmainfont{Latin Modern Roman}
\usepackage{amsmath,amssymb,mathtools,unicode-math}
\setmathfont{Latin Modern Math}
\usepackage{enumitem,needspace,xurl,hyperref}
\setlength{\parindent}{0pt}
\setlength{\parskip}{5pt}
\setlength{\emergencystretch}{3em}
\newcommand{\N}{\mathbb{N}}
\newcommand{\Z}{\mathbb{Z}}
\newcommand{\Q}{\mathbb{Q}}
\newcommand{\R}{\mathbb{R}}
\newcommand{\C}{\mathbb{C}}
\newcommand{\F}{\mathbb{F}}
\newcommand{\A}{\mathbb{A}}
\newcommand{\PP}{\mathbb{P}}
\newcommand{\E}{\mathbb{E}}
\newcommand{\eps}{\varepsilon}
\newcommand{\dd}{\,\mathrm{d}}
\begin{document}
\section*{UnsolvedMath ID: 30005620}
\textbf{Source problem code:} catalogue entry
\subsection{Definitions and mathematical statement}
Let \(M\) be a closed smooth surface, meaning a compact smooth manifold of dimension two without boundary. For a smooth Riemannian metric \(g\) and a real-valued function \(f\in C^\infty(M)\), put \(\sigma=e^{-f}\) and \(\dd\mu_{g,f}=e^{-f}\dd v_g\), where \(\dd v_g\) is the Riemannian area measure. The nonnegative weighted Laplacian specified by the supplied integration identity is
\[
L_{g,f}u=-e^f\operatorname{div}_g(e^{-f}\nabla_g u)
=-\operatorname{div}_g(\nabla_g u)+\langle\nabla_g f,\nabla_g u\rangle_g.
\]
Equivalently, for all \(\phi,\psi\in C^\infty(M)\),
\[
\int_M\langle\nabla_g\phi,\nabla_g\psi\rangle_g e^{-f}\dd v_g
=\int_M\phi L_{g,f}\psi\,e^{-f}\dd v_g.
\]
Its eigenvalues are those of its self-adjoint realization in \(L^2(M,\mu_{g,f})\), counted with multiplicity. One may label the full spectrum by
\[
0=\lambda_0(g,f)\leq\lambda_1(g,f)\leq\lambda_2(g,f)\leq\cdots;
\]
this is an explicit indexing convention for the present definitions, since the supplied record does not specify its indexing convention.

\textbf{Requested bound, with unresolved normalization.} The supplied statement asks for an upper bound on quantities denoted by \(\overline\lambda_i(M,g,e^{-f})\), described as normalized eigenvalues, independent of both \(g\) and \(f\). Once the normalization and index convention have been specified, a fixed-surface formulation of the requested bound would be: for each permitted index \(i\), does there exist a finite constant \(C(M,i)\) such that
\[
\overline\lambda_i(M,g,e^{-f})\leq C(M,i)
\quad\text{for every smooth Riemannian metric }g\text{ on }M
\quad\text{and every }f\in C^\infty(M;\R)?
\]
This display makes explicit one possible dependence on the fixed surface and the index; the record does not determine whether a stronger uniformity over surfaces or indices is intended. No curvature, geometric, or weight restriction is supplied beyond the stated setting.

The record does not define \(\overline\lambda_i\). In particular, it does not specify whether normalization uses Riemannian area, total weighted measure, or another quantity, or whether a normalization of \(f\) is imposed. Consequently, the display is a conditional formulation of the requested type of estimate, not a fully determined mathematical conjecture. The original excerpt also introduces the Bakry--\'Emery question as a similar question after an earlier normalized-eigenvalue question, so it does not itself establish that the earlier notation refers to this operator.
\subsection{Short English statement}
The supplied problem asks whether normalized eigenvalues of the Bakry--\'Emery Laplacian on a closed surface admit upper bounds independent of the metric and every smooth weight. The operator is specified by a weighted integration identity, but the eigenvalue normalization and the permitted dependence of the bound are missing. These must be resolved before the question has a unique precise formulation.
\subsection{Sources}
\begin{itemize}
\item Supplied local UnsolvedMath registry record, permanent ID \(30005620\), exact source problem code \texttt{catalogue entry}, title ``Metric-Independent Bounds for Bakry--\'Emery Eigenvalues.'' The record supplies proposed year 2023 and names no proposer.
\item Supplied report citation: \emph{Geometric Spectral Theory (2024)}, Oberwolfach Reports, DOI \texttt{10.4171/owr/2023/36}, pp.\ 2017--2106. Original report link: \url{https://doi.org/10.4171/owr/2023/36}. The supplied background describes Oberwolfach Reports as records of mathematical workshops and their open-problem discussions.
\item Supplied literature reference: Funano--Shioya, described in the record as ``weighted eigenvalue bounds under curvature'': \url{https://doi.org/10.1007/s00039-013-0215-x}.
\item Literature assessment dated exactly \texttt{2026-08-22}: the supplied curation labels the problem open and states that universal upper bounds are known under additional curvature or geometric hypotheses, while independence from both the metric and arbitrary weight remains unresolved. This is supplied dated curation derived from the Oberwolfach curation, not an independently verified literature assessment; its status and cited claims require independent verification.
\item The record labels the statement \texttt{corrected\_verified} and says it was reconstructed and checked against the official report. These are supplied verification claims, not verification performed for this draft.
\end{itemize}
\end{document}
